\chapter{The Finished Receipt}

The volume has now done two different things. First, it built the instrument: a ladder of distinctions, counts,
faces, traces, gauges, residues, and costs. Then it used the instrument: first as an external recoil manual for
metrological \(\alpha\), then as a Cooper-pair interface where an electron can be carried to a second electron
and then to the superconducting pair boundary. The last question is whether the apparatus can be left in that
state and asked for a finite study of its own reading.

It can. The companion file is
\[
  \texttt{device/Measurement/Meanwhile50.lean}.
\]
That file does not restart the machine. It imports the electron and pairing state already constructed, keeps the
electron state alive, advances the carried orbit and PRNG state one step per trial, and reads the electron-charge
aperture at each step. Each trial is therefore a small experiment:
\[
  \alpha_i = \frac{e_i^2}{4\pi},
  \qquad c=1.
\]
The value \(e_i^2\) is not an arbitrary decimal. The selected face chooses one of the three rational
charge-square readings in the adjacent aperture:
\[
\begin{array}{c|c}
\hbox{selected face} & \hbox{charge-square reading}\\
\hline
\hbox{charge} & \hbox{lower square}\\
\hbox{mass} & \hbox{target square}\\
\hbox{value} & \hbox{upper square}
\end{array}
\]
The carried state decides which face is selected. The division by \(4\pi\) is the normalization already fixed in
the charge computation. In this normalization, \(h=c=1\), and the charge relation reads
\[
  e^2 = 4\pi\alpha.
\]
So the recovery step is the reverse:
\[
  \alpha = \frac{e^2}{4\pi}.
\]

\section{How to Run the Experiment}

Run the experiment from the Lean project directory:
\begin{quote}\small
\begin{verbatim}
cd device
lake build Measurement.Meanwhile50
\end{verbatim}
\end{quote}
This is the safest command because it builds the module and every dependency it needs. A reader who wants to run
only the file after the project is already compiled may also use:
\begin{quote}\small
\begin{verbatim}
lake env lean Measurement/Meanwhile50.lean
\end{verbatim}
\end{quote}
The file prints three reports that matter for this chapter:
\begin{quote}\small
\begin{verbatim}
#eval defaultElectronAlphaStudyReport?
#eval defaultElectronAlphaPathRelationReport?
#eval defaultElectronAlphaRichardsonReport?
\end{verbatim}
\end{quote}
Other imported modules may also print receipts when the build recompiles dependencies. Those are useful context,
but the three names above are the experiment's final readout. The first is the one-hundred-trial alpha study. The
second is the comparison of the path \(1\to2\to3\) against the direct path \(1\to3\). The third is the final
Richardson check: three finite alpha studies, two signed refinement slips, and a bracketed extrapolated alpha
reading.

The run has four preparation steps, all visible in the imported modules:
\begin{enumerate}
\item \(\texttt{Meanwhile48}\) plugs the turn decider into the corridor.
\item \(\texttt{Meanwhile49}\) builds \(\texttt{Fact.ELECTRON}\), creates the first electron readings, and
recovers \(\alpha\) from \(e^2/(4\pi)\).
\item \(\texttt{Meanwhile50}\) advances the carried orbit/PRNG state one trial at a time.
\item Each trial chooses charge, mass, or value and maps that face to the lower, target, or upper charge-square
reading in the adjacent aperture.
\end{enumerate}
The machine is therefore not restarted between trials. The seed and residue in trial \(i+1\) are the output of
trial \(i\).

To change the number of trials, make a local report with a different final argument:
\begin{quote}\small
\begin{verbatim}
def myAlphaStudy? : Option ElectronAlphaStudyReport :=
  defaultAdjacentWobbleReport?.map fun aperture =>
    electronAlphaStudyReport aperture
      cavendishChargeMassNormalization 1000

#eval myAlphaStudy?
\end{verbatim}
\end{quote}
The last number is the trial count. The report shape does not change when the count changes; only the sample
list, face counts, mean, variance, seed, and final residue change.

\section{The Study Object}

The word ``study'' is not decoration. The early apparatus defined a study as a chain of data under a fact:
\[
  \texttt{Study.hypothesis}
  \quad\hbox{or}\quad
  \texttt{Study.data}.
\]
The alpha run keeps that shape. The formal side records one hundred data entries under
\(\texttt{Fact.ELECTRON}\). The numeric side records the sample list and its moments. That separation matters.
The study is the proof-shaped ledger saying that one hundred trials were performed; the statistics are the
number-shaped ledger saying what those trials read.

The report is named:
\begin{quote}\small
\begin{verbatim}
defaultElectronAlphaStudyReport?
\end{verbatim}
\end{quote}
and the finished receipt begins:
\begin{quote}\small
\begin{verbatim}
some { name := "electron-alpha-carried-wobble-study",
  normalization := "h = c = 1",
  cSetToOne := true,
  startsFromElectronState := true,
  startingCompletedTurns := 2,
  trialCount := 100,
  formalStudyDepth := 100,
  ...
}
\end{verbatim}
\end{quote}
The study starts after the electron has completed two turns. That is the state reached by the electron fact in
the prior file. The run then advances exactly one hundred trials.

The essential fields are:
\begin{itemize}
\item \(\texttt{trialCount}\): how many carried trials were run.
\item \(\texttt{formalStudyDepth}\): how many \(\texttt{Study.data}\) entries were recorded.
\item \(\texttt{initialSeed}\) and \(\texttt{finalSeed}\): the PRNG state before and after the study.
\item \(\texttt{initialResidueNumerator}\) and \(\texttt{finalResidueNumerator}\): the carried orbit residue
before and after the study.
\item \(\texttt{chargeFaceSamples}\), \(\texttt{massFaceSamples}\), and \(\texttt{valueFaceSamples}\): how
often the decider chose each aperture face.
\item \(\texttt{sampleAlphaMeanScaledAt18}\): the mean of the sampled alpha readings.
\item \(\texttt{sampleAlphaVarianceScaledAt36}\): the variance before dropping back to \(10^{18}\) scale.
\end{itemize}
The equality \(\texttt{trialCount}=\texttt{formalStudyDepth}\) is the first sanity check. It says the numeric
study and the formal study chain counted the same number of trials.

\section{The Hundred Readings}

The hundred selections are nearly perfectly balanced:
\[
\begin{array}{c|c}
\hbox{face} & \hbox{number of trials}\\
\hline
\hbox{charge} & 33\\
\hbox{mass} & 34\\
\hbox{value} & 33
\end{array}
\]
The extra mass reading is not hand-chosen. It is the finite residue of one hundred draws across three faces.
The target, lower, and upper alpha values read:
\[
\begin{array}{c|c}
\hbox{face} & \alpha\ \hbox{read at }10^{-18}\hbox{ scale}\\
\hline
\hbox{charge} & 0.007298662680727034\\
\hbox{mass} & 0.007298668569525773\\
\hbox{value} & 0.007298710880818371
\end{array}
\]
The source target inside the aperture is therefore
\[
  \alpha_{\rm source}
  =
  0.007298668569525773,
\]
with
\[
  \alpha_{\rm source}^{-1}
  =
  137.011290548979455469.
\]

The one-hundred-trial study returns:
\[
\begin{array}{l|r}
\hbox{source alpha, scaled by }10^{18} & 7298668569525773\\
\hbox{sample mean alpha, scaled by }10^{18} & 7298680588948746\\
\hbox{absolute mean-source difference, scaled by }10^{18} & 12019422973\\
\hbox{sample mean inverse alpha, scaled by }10^{18} & 137011064919616304245\\
\hbox{sample variance, scaled by }10^{18} & 457\\
\hbox{sample variance, scaled by }10^{36} & 457758203912487384985
\end{array}
\]
In ordinary decimal notation:
\[
\begin{aligned}
  \alpha_{\rm mean} &\approx 0.007298680588948746,\\
  \alpha_{\rm mean}^{-1} &\approx 137.011064919616304245,\\
  {\rm Var}(\alpha) &\approx 4.57758203912487384985\times 10^{-16}.
\end{aligned}
\]
The mean sits above the source because the finite one-hundred-trial run selected the upper face and lower face
equally often but selected the target face one extra time. The variance is not laboratory uncertainty; it is the
spread of this internal aperture study. It says how wide the device reading is when the carried decider samples
the adjacent rational neighborhood.

\section{Richardson Extrapolation}

The last proof-facing step is not another physical assumption. It is a finite refinement check on the alpha
study itself. The cheap receipt used \(100,200,400\) trials while the mechanism was being built. There is no
reason to stop there. The finished receipt runs the same carried experiment at three larger depths:
\[
\begin{array}{c|c|c|c|c}
\hbox{trials} & \hbox{charge} & \hbox{mass} & \hbox{value}
  & \bar\alpha\hbox{ at }10^{-18}\hbox{ scale}\\
\hline
10000 & 3389 & 3337 & 3274 & 7298680426529076\\
20000 & 6621 & 6655 & 6724 & 7298680845095521\\
40000 & 13349 & 13332 & 13319 & 7298680692889066
\end{array}
\]
These are not three unrelated reruns. They are nested finite studies of the same carried machine: the same initial
electron state, the same aperture, the same PRNG law, and a longer allowance for the state to keep moving.

The two refinement slips are the signed changes in the mean:
\[
\begin{aligned}
  \delta_{10000\to20000} &= +0.000000000418566445,\\
  \delta_{20000\to40000} &= -0.000000000152206455.
\end{aligned}
\]
The sign change is the important event. The refinement passed through a turn of the finite sampling error, so the
last read should not be a single naked mean. It should be bracketed.

For a first-order finite-depth error, the Richardson projection from a pair of studies is
\[
  R_{n,2n}=2\bar\alpha_{2n}-\bar\alpha_n.
\]
The two projections are therefore
\[
\begin{aligned}
  R_{10000,20000} &= 0.007298681263661966,\\
  R_{20000,40000} &= 0.007298680540682611.
\end{aligned}
\]
They give the final bracket
\[
  0.007298680540682611
  \leq
  \alpha_{\rm Richardson}
  \leq
  0.007298681263661966,
\]
with width
\[
  0.000000000722979355.
\]
The midpoint receipt is
\[
  \alpha_{\rm Richardson}
  \approx
  0.007298680902172288.
\]
Lean records the two checks that make this a bracket rather than decoration:
\begin{quote}\small
\begin{verbatim}
signedSlipsBracketZero := true
fineStudyMeanInsideRichardsonBracket := true
\end{verbatim}
\end{quote}
The same report also carries the two orbit-slip brackets from the corridor. Those are the angular places where
the apparatus counted turns. The Richardson slips are the numerical echoes of those turn cuts in the alpha mean.
That is the last step of the proof: three studies, two slips, and the extrapolated alpha held between them.

This larger run also clarifies the interpretation. The high-depth mean is not trying to collapse back to the
mass-face source alone. It is reading the carried decider's long-run aperture: charge, mass, and value sampled as
faces of the same corridor. The source alpha remains the target square inside the aperture; the Richardson receipt
is the apparatus reading the aperture distribution after many carried trials.

\section{The Path Relation}

The last check compares the two routes through the electron corridor:
\[
  1\to2\to3
  \qquad\hbox{and}\qquad
  1\to3.
\]
The corridor residues compose exactly:
\[
\begin{array}{l|r}
1\to2 & 1\\
2\to3 & 1\\
1\to2\to3 & 2\\
1\to3 & 2
\end{array}
\]
Lean records this as
\begin{quote}\small
\begin{verbatim}
corridorResiduesRelated := true
\end{verbatim}
\end{quote}
So the corridor says that going from \(1\) to \(2\) and then \(2\) to \(3\) is the same residue as going directly
from \(1\) to \(3\).

The alpha aperture adds the finer reading. The three rungs are:
\[
\begin{array}{c|c|c}
\hbox{rung} & \hbox{face} & \alpha\\
\hline
1 & \hbox{mass} & 0.007298668569525773\\
2 & \hbox{value} & 0.007298710880818371\\
3 & \hbox{charge} & 0.007298662680727034
\end{array}
\]
The signed differences compose:
\[
\begin{aligned}
  \Delta_{1\to2} &= +0.000000042311292598,\\
  \Delta_{2\to3} &= -0.000000048200091337,\\
  \Delta_{1\to2\to3} &= -0.000000005888798739,\\
  \Delta_{1\to3} &= -0.000000005888798739.
\end{aligned}
\]
Lean records:
\begin{quote}\small
\begin{verbatim}
signedDeltasRelated := true
\end{verbatim}
\end{quote}
But the absolute path lengths do not agree:
\[
\begin{aligned}
  |\Delta_{1\to2}|+|\Delta_{2\to3}|
    &= 0.000000090511383935,\\
  |\Delta_{1\to3}|
    &= 0.000000005888798739.
\end{aligned}
\]
The triangle excess is
\[
  0.000000084622585196.
\]
Lean records:
\begin{quote}\small
\begin{verbatim}
absolutePathLengthsRelated := false
\end{verbatim}
\end{quote}

That is the important shape. The middle rung is not invisible. It cancels in the signed displacement, so the
direct and composed routes land at the same signed alpha difference. But it does not cancel as cost. The route
through \(2\) is a real excursion through the aperture. The device can therefore distinguish equality of endpoint
displacement from equality of path length. That distinction is exactly what a measuring instrument must be able
to keep.

\section{What the Experiment Means}

The experiment is an internal apparatus study, not a CODATA replacement. It does not say that the physical world
has been measured to CODATA precision by a laptop build. It says something more local and more important for this
volume: once the construction has produced an electron fact, a corridor, a carried decider, and an
\(e^2/(4\pi)\) recovery rule, the machine can be left in state and can produce a repeatable finite distribution
of alpha readings.

The source alpha is the target square inside the adjacent charge aperture. The mean alpha is what the carried
decider reads after one hundred trials. The variance is the spread of the sampled aperture, not an experimental
standard uncertainty from laboratory hardware. To turn this into metrology, attach the recoil route from the
previous part, supply real laboratory tapes and correction ledgers, and compare against external route-by-route
measurements. To study the apparatus itself, change the trial count and watch the mean, variance, face counts,
seed, and residue move.

The path relation gives the geometric lesson. The corridor residue is additive:
\[
  (1\to2)+(2\to3) = (1\to3).
\]
The signed alpha displacement is also additive. But the absolute path length is not. The middle rung is a
detour. That is why the apparatus keeps both signs and costs: two routes can agree as endpoint displacement
while disagreeing as work spent along the way.

Read the final receipt in that order:
\begin{enumerate}
\item the run was finite and budgeted;
\item the state was carried from trial to trial;
\item the face counts show what the decider sampled;
\item the mean and variance summarize the aperture readings;
\item the Richardson report uses three studies and two signed slips to bracket the extrapolated alpha;
\item the path relation separates endpoint equality from path cost.
\end{enumerate}
That is the completed alpha receipt. The only remaining question is how the simulation that produced the receipt
was itself measured.

\section{What the Receipt Hands Forward}

The apparatus has:
\begin{enumerate}
\item built the electron fact;
\item made the electron distinguishable, admissible, and countable;
\item created the first three electron readings by number;
\item compared the quartet and found the first repeated face;
\item recovered \(\alpha\) from \(e^2/(4\pi)\) in the \(h=c=1\) normalization;
\item left the machine state alive and performed a one-hundred-trial alpha study;
\item computed exact scaled-Nat readings for the study mean and variance;
\item ran the three-study Richardson check and bracketed the extrapolated alpha;
\item compared \(1\to2\to3\) with \(1\to3\) and separated signed relation from path cost.
\end{enumerate}

This is the end of the construction proper because the book has reached the thing it promised: a finite machine
that can make a reading and hand the reader the exact place where interpretation enters. The recoil chapters show
how to attach the machine to an external metrology route. The pairing chapters show how to add the
electron-electron-Cooper-pair interface and stop an otherwise endless computation by cost. This study shows the
machine reading its own alpha aperture while carrying state from trial to trial.

Nothing in that receipt asks the reader to worship the digits. The digits are not the point. The point is that
the route by which the digits appeared is now on the page: typed, budgeted, graded, and repeatable.

The last act now measures the simulation itself. A reader can run the Lean file, see the same receipt, change the
budget, change the trial count, change the aperture, or replace the demonstration scales with laboratory scales.
The apparatus will make the change visible. The receipt is no longer only a number. It is a number, a route, and
an elaborator cost waiting to be calibrated.
